\documentclass[a4paper,12pt]{article}
\usepackage[T2A]{fontenc}
\usepackage[utf8]{inputenc}
\usepackage[russian]{babel}
\usepackage{amsmath,amssymb,mathtools}
\usepackage{geometry}
\geometry{left=22mm,right=20mm,top=20mm,bottom=20mm}
\usepackage{microtype}
\usepackage{array,tabularx,booktabs}
\usepackage{xcolor}
\usepackage{tikz}
\usepackage{fancyhdr}
\usepackage{hyperref}
\definecolor{accent}{HTML}{087C83}
\definecolor{ink}{HTML}{252A31}
\definecolor{muted}{HTML}{626D76}
\hypersetup{colorlinks=true,linkcolor=accent,pdfauthor={Лёвин Артём Александрович},pdftitle={Ревью домашней работы — 29 октября 2026 г.}}
\setlength{\parindent}{0pt}
\setlength{\parskip}{4pt}
\linespread{1.07}
\setlength{\headheight}{14pt}
\pagestyle{fancy}
\fancyhf{}
\fancyhead[L]{\footnotesize\color{muted}МАТЕМАТИКА \quad / \quad Ревью домашней работы}
\fancyhead[R]{\footnotesize\color{muted}29.10.2026}
\fancyfoot[C]{\footnotesize\color{muted}Лёвин Артём Александрович эксклюзивно для Николь Саркисьянц}
\renewcommand{\headrulewidth}{0pt}
\renewcommand{\footrulewidth}{0pt}
\newcommand{\smalllabel}[1]{\par\smallskip\textbf{#1}\par\nobreak}
\newcommand{\errlabel}{\textcolor{red!75!black}{\textbf{Ошибка:}}}
\newcommand{\ideallabel}{\textcolor{blue!60!black}{\textbf{Ключевая идея:}}}
\newcommand{\finalmsg}{\textcolor{accent}{\textbf{Итог:}}}
\begin{document}
\thispagestyle{empty}
\begin{center}
\vspace*{39mm}
{\color{accent}\large МАТЕМАТИКА}\par
\vspace{22mm}
{\Huge Ревью домашней работы}\par
\vspace{12mm}
{\color{accent}\Large Показательные и рациональные неравенства}\par
\vspace{33mm}
{\large Лёвин Артём Александрович}\par
\vspace{13mm}
{\large 29 октября 2026 г.}\par
\vfill
{\small\color{muted}Лёвин Артём Александрович}\par
{\small\color{muted}эксклюзивно для Николь Саркисьянц}\par
\vspace{19mm}
\begin{tikzpicture}
\draw[accent,line width=0.72pt] (0,0) .. controls (1.15,0.42) and (2.35,-0.35) .. (3.45,0)
.. controls (4.55,0.33) and (5.35,-0.28) .. (6.50,0);
\end{tikzpicture}
\end{center}
\clearpage
\section*{Сводная таблица}
\vspace{4mm}
В таблице приведено задание, в решении которого обнаружена ошибка.

\vspace{4mm}\par\noindent
\small
\renewcommand{\arraystretch}{1.5}
\begin{tabularx}{0.98\linewidth}{@{}p{8mm}p{27mm}X p{35mm}p{22mm}@{}}
\toprule
\textbf{№} & \textbf{Тема} & \textbf{Суть ошибки} & \textbf{Ваш ответ} & \textbf{Верный ответ}\\
\midrule
\hyperref[task:9]{9} & Рациональное неравенство & При приведении подобных слагаемых потерян знак минус: $2x-3x=-x$. & $\begin{gathered}(-\infty;-13]\\ \cup(4;+\infty)\end{gathered}$ & $(4;13]$\\
\bottomrule
\end{tabularx}
\normalsize
\vfill
{\color{muted}\small Нажмите на номер задания, чтобы перейти к подробному разбору.}
\clearpage
\thispagestyle{empty}
\vspace*{0.38\textheight}
\begin{center}
{\color{accent}\huge Разбор ошибок}
\end{center}
\clearpage
\phantomsection\label{task:9}
\section*{Задание № 9}
\smalllabel{Условие}
Решите неравенство:
\[
\frac{2x+1}{x-4}\ge 3.
\]
\smalllabel{Ваше решение}
Правильно перенесли тройку в левую часть и привели дроби к общему знаменателю:
\[
\frac{2x+1}{x-4}-3\ge0,
\qquad
\frac{2x+1-3x+12}{x-4}\ge0.
\]
В следующей строке записано
\[
\frac{x+13}{x-4}\ge0.
\]
По этой дроби выполнен метод интервалов.

\smalllabel{Ваш ответ}
\[
 x\in(-\infty;-13]\cup(4;+\infty).
\]
\smalllabel{Ошибка}
\errlabel\ Неправильно приведены подобные слагаемые в числителе.

\smalllabel{Где именно возникает ошибка}
\textbf{Последний верный шаг:}
\[
\frac{2x+1-3x+12}{x-4}\ge0.
\]
\textbf{Первый неверный шаг:}
\[
\frac{\color{red!75!black}{x}+13}{x-4}\ge0.
\]
Действительно, $2x-3x=(2-3)x=-x$, а $1+12=13$.
Поэтому в числителе должно стоять $13-x$.

\smalllabel{Почему это неверно}
Из-за потерянного минуса Вы стали исследовать другую дробь.
Числитель $x+13$ обращается в нуль при $x=-13$,
а настоящий числитель $13-x$ обращается в нуль при $x=13$.
Эти выражения имеют разные знаки на числовой прямой,
следовательно, полученные интервалы не могут совпасть.

\smalllabel{Ключевая идея}
\ideallabel\ При вычитании выражений сначала раскройте скобки,
затем сложите коэффициенты при переменной и свободные члены отдельно.
Для дробного неравенства определите нули числителя и знаменателя,
а знак дроби устанавливайте по обоим множителям.
\clearpage
\section*{Задание № 9. Продолжение}
\smalllabel{Исправление от последнего верного шага}
\[
\begin{aligned}
\frac{2x+1-3x+12}{x-4}&\ge0,\\[3pt]
\frac{(2-3)x+(1+12)}{x-4}&\ge0,\\[3pt]
\frac{13-x}{x-4}&\ge0.
\end{aligned}
\]
\smalllabel{Полное правильное решение}
\textbf{Шаг 1.} Сразу отметим ограничение: $x-4\ne0$, то есть $x\ne4$.

\textbf{Шаг 2.} Найдём точку, где числитель равен нулю:
\[
13-x=0\quad\Longleftrightarrow\quad x=13.
\]
Точка $13$ может входить в ответ, потому что знак неравенства нестрогий.

\textbf{Шаг 3.} Знаменатель обращается в нуль при $x=4$.
Точку $4$ исключаем: делить на нуль нельзя.

\textbf{Шаг 4.} Расставим знаки на трёх промежутках:
\[
\renewcommand{\arraystretch}{1.35}
\begin{array}{c|ccc}
 & x<4 & 4<x<13 & x>13\\ \hline
 13-x & + & + & -\\
 x-4  & - & + & +\\ \hline
 \dfrac{13-x}{x-4} & - & + & -
\end{array}
\]
Дробь положительна между $4$ и $13$. В точке $13$ она
равна нулю, что допускается знаком $\ge$.

\textbf{Шаг 5.} Запишем найденный промежуток:
\[
\boxed{x\in(4;13]}.
\]
\smalllabel{Проверка}
При $x=5$ исходная дробь равна $11$; неравенство $11\ge3$ верно.
При $x=13$ получаем $27/9=3$, поэтому точка $13$ включается.
При $x=14$ имеем $29/10<3$, значит правее $13$ решений нет.
Число $4$ подставлять нельзя: знаменатель равен нулю.

\smalllabel{Итог}
\finalmsg\ Правильный ответ: $\boxed{(4;13]}$.
Критический навык --- аккуратно приводить подобные слагаемые
перед применением метода интервалов.

\smalllabel{Задача на закрепление}
Решите неравенство:
\[
\frac{3x+2}{x-5}\ge4.
\]
\end{document}
